\documentclass[a4paper,10pt]{scrartcl}

\usepackage[ngerman]{babel}
\usepackage[utf8]{inputenc}
\usepackage[T1]{fontenc}
\usepackage[babel,german=quotes]{csquotes}
\usepackage{hyperref}
\usepackage{geometry}
\usepackage{graphicx}
\usepackage{enumitem}
\usepackage{multicol}

\setlength\parindent{0pt}
\setitemize{itemsep=0pt}


\geometry{
	a4paper,
	left=20mm,
	top=20mm,
	bottom=15mm,
	footskip=5mm
}

\subject{Gruppe \textit{Nummer}}
\title{Entwicklerhandbuch für das Thema \glqq\textit{Titel}\grqq}
\subtitle{Betreuer: \textit{Name}}
\author{\textit{Namen der Teilnehmenden}}
\date{\textit{Datum}}

\begin{document}

\maketitle

\textit{\textbf{Anmerkung:} Alles, was in diesem Dokument kursiv geschrieben ist, muss von der Gruppe entsprechend ihres Themas angepasst (und diese und alle übrigen Anmerkungen entfernt) werden}

\tableofcontents

\newpage

\section{Beschreibung}

\textit{Kurze, textuelle Beschreibung der entwickelten Software, ihrer Bestandteile und ihrer wichtigsten Funktionen}

\newpage

\section{Frontend / Oberfläche}

\textit{\textbf{Anmerkung:} Sofern die Anwendung kein Frontend bzw. keine separate Oberfläche besitzt, kann dieser Abschnitt weggelassen werden}

\subsection{Übersicht}

\textit{Grobe Beschreibung des Aufbaus und der wichtigsten Komponenten des Frontends}

\subsection{Datei-/Ordnerstruktur}

\textit{Stichpunktartige Erläuterung der Datei- und Ordnerstruktur des Frontend-Teils des Projekts; dabei muss nicht jede einzelne Datei beschrieben werden, sondern eher der grobe Aufbau (also z.B. was sich in welchem Ordner befindet)}

\subsection{Genutzte Abhängigkeiten}

\begin{itemize}
    \item \textit{Auflistung und jeweils kurze Erläuterung der verwendeten Tools/Frameworks/Bibliotheken}
    \item \textit{...}
\end{itemize}

\subsection{Einrichten der Entwicklungsumgebung}

\textit{Erläuterung der Schritte (und ggf. dafür zu installierende Software) um das Projekt für die lokale Entwicklung einzurichten}

\subsection{Build}

\textit{Erläuterung der Schritte zum Bauen des Frontends (sofern möglich explizit die Erläuterung eines Bauens für einen Produktiveinsatz, während ein Bauen für den Einsatz während der Entwicklung im vorherigen Abschnitt erläutert werden soll)}

\newpage

\section{Backend}

\textit{\textbf{Anmerkung:} Sofern die Anwendung kein Backend besitzt, kann dieser Abschnitt weggelassen werden; falls ein Backend vorhanden ist, dann sollen die Abschnitte jeweils vom Inhalt her analog zu den Erläuterungen für den Frontend-Abschnitt gefüllt werden}

\subsection{Übersicht}

\textit{...}

\subsection{Datei-/Ordnerstruktur}

\textit{...}

\subsection{Genutzte Abhängigkeiten}

\begin{itemize}
    \item \textit{...}
\end{itemize}

\subsection{Einrichten der Entwicklungsumgebung}

\textit{...}

\subsection{Build}

\textit{...}

\newpage

\section{Bauen und Starten des gesamten Projekts}

\textit{\textbf{Anmerkung:} Sofern die Anwendung aus mehreren Bestandteilen (z.B. Front- und Backend) besteht, dann sollen in diesem Abschnitt stichpunktartig die Schritte zum Bauen und Starten des kompletten Projekts erläutert werden}

\end{document}